\chapter{Introduction}
\label{ch:introduction}

\section{Origin of the Report}
The Bachelor of Science in Software Engineering programme offered by the
Department of Software Engineering of the Institute of Information and
Communication Technology (IICT), Shahjalal University of Science and
Technology (SUST), requires every final-year student to complete a
six-month industry internship under the course code \textbf{SWE--420}.
The course is an eighteen-credit lab course that bridges the gap between
classroom learning and real-world software engineering practice,
exposing students to the technologies, processes, teams, and culture
they will encounter as professionals.

This report has been prepared as the formal academic deliverable of that
internship. It documents my experience at \textbf{FlowGenX AI}, a
Silicon Valley-based software development company building an agentic
integration platform that combines multi-agent orchestration, low-code
workflow design, and a large catalogue of native connectors. The
document was prepared under the guidance of my academic department and
the company's engineering leadership.

\section{Objectives of the Report}
The internship programme, and by extension this report, has both an
academic and a personal objective. The specific goals I set for myself
during the internship and that this report aims to demonstrate are:

\begin{enumerate}[leftmargin=*]
    \item \textbf{Bridging theory and practice.} To translate the
    theoretical concepts of software engineering acquired across seven
    semesters --- programming languages, databases, software
    architecture, software engineering principles --- into hands-on,
    production-grade contributions on a real product.

    \item \textbf{Learning industry-standard engineering practices.}
    To work inside a disciplined feature-branch Git workflow, write
    code that passes peer review, build and run automated test suites,
    and contribute to a continuous integration pipeline used by paying
    customers.

    \item \textbf{Working on a meaningful, production-grade project.}
    To take complete ownership of features --- in my case, end-to-end
    connector development --- from API research and design through
    implementation, testing, documentation, and live deployment.

    \item \textbf{Developing depth in modern Generative AI engineering.}
    To gain practical exposure to large language model providers, the
    Model Context Protocol (MCP), Agent-to-Agent (A2A) communication,
    Retrieval-Augmented Generation (RAG) knowledge bases, and the
    workflow engines that compose these primitives into real
    applications.

    \item \textbf{Building professional non-technical skills.} To grow
    as a remote-first engineer collaborating across time zones with a
    globally distributed team, communicating clearly through text and
    video, planning sprint work, raising blockers early, and giving
    constructive code review to peers.

    \item \textbf{Understanding modern startup engineering culture.}
    To experience first-hand how a small, ambitious team builds a
    serious product that competes with established platforms in the
    integration space.

    \item \textbf{Documenting the experience faithfully.} To produce a
    report that honestly reflects the work, the learnings, and the
    challenges of the internship --- so that it serves as both an
    academic deliverable and a record I can return to in my own future.
\end{enumerate}

\section{Scope of the Report}
This report describes the company, the engineering process, the
projects I worked on, and the skills I acquired during my internship at
FlowGenX AI between November 2025 and May 2026. Specifically, the
report covers:

\begin{itemize}[leftmargin=*]
    \item An overview of the internship programme as it was structured
    at FlowGenX AI: the team I was placed on, the supervisor and
    teammates I worked with, the project I contributed to, and the
    duration and weekly cadence of the work.

    \item A profile of FlowGenX AI as an organisation --- its products,
    its agentic platform architecture, its target market, the
    technologies it uses, the office environment (including the
    realities of remote work), and the company culture I experienced.

    \item The software engineering methodology and process I followed
    at the company: the development lifecycle, branching strategy,
    code review and testing practices, the connector development
    pipeline, and the tools used across version control, project
    management, code quality, and deployment.

    \item A detailed account of my actual contributions, organised in
    two phases --- an initial onboarding and learning period followed
    by sustained live project work on FlowGenX's
    \texttt{runtime\_connectivity} project. This includes specific
    connectors I designed and shipped, bugs I fixed, refactoring I
    drove, and documentation I authored.

    \item The technical and non-technical skills I acquired during the
    internship and the personal reflections that came out of working
    inside a remote-first global startup.
\end{itemize}

The report focuses only on what falls within the scope of an academic
deliverable. It is not a technical specification of FlowGenX's product,
nor an exhaustive engineering write-up of every commit; it is a faithful
reflection of an intern's six-month journey inside the company.

\section{Methodology}
The information presented in this report has been compiled from
multiple complementary sources to ensure both accuracy and personal
authenticity:

\begin{itemize}[leftmargin=*]
    \item \textbf{Direct work experience.} The bulk of this report is
    drawn directly from my five and a half months of full-time work
    inside the company --- the projects I shipped, the meetings I
    attended, the people I worked with, and the artefacts I produced.

    \item \textbf{Version control history.} For Chapter~4 in particular,
    I systematically reviewed my Git commit history on the
    \texttt{runtime\_connectivity} repository to compile an accurate
    timeline of what was delivered when, including branch names,
    merged pull requests, and per-period focus areas.

    \item \textbf{Internal repository documentation.} I referred to the
    project's \texttt{README.md} and \texttt{CLAUDE.md} files (the
    latter being the team's onboarding-and-conventions document for
    AI-assisted code editors) to describe the engineering process,
    architecture, and conventions accurately.

    \item \textbf{Public-facing company materials.} Company background,
    vision, product positioning, platform pillars, and use cases were
    drawn from the company's public website
    (\texttt{https://www.flowgenx.ai/}, including the \emph{About},
    \emph{Platform}, \emph{Use Cases}, and \emph{FAQ} pages).

    \item \textbf{Direct communication with company leadership.} For
    factual gaps that the public website does not address (such as
    employee count and team composition), I have where appropriate
    reached out to the leadership team for confirmation rather than
    speculate.
\end{itemize}

\section{Limitations}
Although this report has been prepared with care, three limitations
should be acknowledged upfront:

\begin{enumerate}[leftmargin=*]
    \item \textbf{Confidentiality and intellectual property
    constraints.} My internship offer letter establishes a strict
    Intellectual Property and Confidentiality clause. As a result,
    proprietary internal architecture, business logic, customer data,
    pricing details, and unreleased product roadmap items have been
    deliberately excluded from this report. Wherever such details are
    relevant, only the publicly disclosed or genuinely non-sensitive
    aspects are described.

    \item \textbf{Personal vantage point.} The narrative is written
    from my own perspective as an intern on a single team --- the
    connectivity team --- working on a single project
    (\texttt{runtime\_connectivity}). FlowGenX AI is a multi-team
    organisation building a multi-product platform; my view of the
    rest of the company is therefore second-hand and necessarily
    incomplete.

    \item \textbf{Remote-first nature of the engagement.} Because the
    internship was fully remote, the chapter on \emph{Life at FlowGenX
    AI} (Chapter~6) does not describe an in-office experience the way
    a co-located internship would. It instead reflects what daily life
    looks like as a remote engineer inside a globally distributed
    team --- an experience that is increasingly common in the modern
    software industry.
\end{enumerate}

Within these honest limitations, this report aims to be a complete and
genuine reflection of my internship at FlowGenX AI.
